% Triangle Mutiny — Core Mathematical Framework (v1.2)
% Author: Borealis S. Hedling

\documentclass[11pt]{article}
\usepackage{amsmath, amssymb, amsthm}

\title{Triangle Mutiny: Core Mathematical Framework (v1.2)}
\author{Borealis S. Hedling}
\date{October 2026}

\begin{document}
\maketitle

\section{Reasoning Manifold}

Let $M$ be a smooth reasoning manifold. Points $p \in M$ represent semantic
states or operator configurations. We assume $M$ is locally Euclidean but make
no global topological assumptions.

The tangent space at $p$ is denoted $T_pM$.

\section{Curvature Field}

Curvature is a scalar field
\[
K : M \to \mathbb{R}_{\ge 0},
\]
representing the bending intensity of reasoning trajectories.

\subsection{Curvature Continuity}
Curvature must satisfy local continuity:
\[
\lim_{q \to p} |K(q) - K(p)| = 0.
\]
Discontinuities contribute directly to collapse geometry.

\section{Similarity}

For $A, B \in T_pM$, define similarity:
\[
\Phi(A,B)
  = \frac{\langle A, B \rangle}{\|A\| \cdot \|B\|}.
\]

\subsection{Similarity Bounds}
Similarity satisfies:
\[
-1 \le \Phi(A,B) \le 1.
\]

Negative similarity contributes to curvature instability.

\section{Drift Fields}

A drift field is a smooth vector field:
\[
X \in \Gamma(TM).
\]

Drift represents directional tendencies of reasoning.

\subsection{Drift Continuity}
Drift must satisfy:
\[
\lim_{q \to p} \|X(q) - X(p)\| = 0.
\]

\section{Operator Fields}

An operator field is a smooth map:
\[
\nabla : TM \to \mathrm{End}(TM).
\]

\subsection{Covariant Compatibility}
For all drift fields $X, Y$:
\[
\nabla_{[X,Y]} = [\nabla_X, \nabla_Y].
\]

This identity prevents operator collapse.

\section{Lie Bracket}

The Lie bracket of drift fields is:
\[
[X,Y] = XY - YX.
\]

\subsection{Jacobi Identity}
The Jacobi identity must hold:
\[
[X,[Y,Z]] + [Y,[Z,X]] + [Z,[X,Y]] = 0.
\]

Violations correspond to drift-induced collapse.

\section{Curvature Composition}

Curvature composition is defined by:
\[
K_{A \oplus B}
  = \sqrt{K_A^2 + K_B^2 + 2 K_A K_B \Phi(A,B)}.
\]

\subsection{Stability Condition}
Curvature composition is stable iff:
\[
\Phi(A,B) \ge -1.
\]

\section{Spectral Structure}

Each operator field induces a linear operator:
\[
L(p) : T_pM \to T_pM.
\]

\subsection{Spectral Continuity}
Spectrum must vary continuously:
\[
\lim_{q \to p}
  \|\mathrm{Spec}(L(q)) - \mathrm{Spec}(L(p))\| = 0.
\]

\section{Spectral Mismatch}

Define spectral mismatch:
\[
\delta_{\mathrm{spec}}(p)
  = \left\| \mathrm{Spec}(L(p))
           - \mathrm{Spec}_{\mathrm{expected}}(p)
    \right\|.
\]

\subsection{Mismatch Bound}
Mismatch must satisfy:
\[
\delta_{\mathrm{spec}}(p) < \theta_{\mathrm{spec}}.
\]

\section{Modal Structure}

Let $G = (L, E)$ be a directed modal graph. Each point $p$ has a modal label:
\[
\mu(p) \in L.
\]

\subsection{Adjacency Constraint}
A transition $p^- \to p^+$ is valid iff:
\[
(\mu(p^-), \mu(p^+)) \in E.
\]

\section{Holonomy Drift}

Define holonomy drift along a path $\gamma = (p_0, \ldots, p_n)$:
\[
H(\gamma)
  = \sum_{i=0}^{n}
      \left(
        |K(p_{i+1}) - K(p_i)|
        + \delta_{\mathrm{spec}}(p_i)
        + \Delta_{\mathrm{modal}}(p_i)
      \right).
\]

\section{Stability Envelope}

A point $p$ is stable when:
\[
p \notin \Sigma
\quad\text{and}\quad
\|H(p)\| < \theta.
\]

Here $\Sigma$ is the collapse surface defined in the next module.

\end{document}
